% GENERATED FILE. Edit canonical lessons, then run npm run generate:books.
% canonical-source-hash: fnv1a64:e3d38fd2397ce81a
% canonical-lessons: GU-C265-prem-karvo, GU-C265-paranvu, GU-C265-sagai, GU-C265-vahu, GU-C265-vishvas

\chapter{Love and marriage}
\label{ch:gu-love-and-marriage}
\hlchaptermodality{265}

\begin{chapteropening}
\textbf{By the end of this chapter:} \emph{I can ask about love, engagements and weddings.}

Complete the last lesson of chapter 265: I can ask about love, engagements and weddings.
\end{chapteropening}

\section[\texorpdfstring{\gu{પ્રેમ} \gu{કરવો} (prem karvo)}{પ્રેમ કરવો (prem karvo)}]{\gu{પ્રેમ} \gu{કરવો} (prem karvo) — to love}

\label{lesson:GU-C265-prem-karvo}



\begin{quote}
Before the new one: say the Gujarati for to recognize, then the Gujarati for to embrace.
\end{quote}

\subsection*{You'll want to know: \gu{પ્રેમ} \gu{કરવો}}
\textbf{\gu{પ્રેમ} \gu{કરવો}} — \emph{prem karvo} — \textquotedblleft{}to love\textquotedblright{}.

\begin{quote}
\textbf{\gu{શું} \gu{તે} \gu{તેની} \gu{પત્નીને} \gu{પ્રેમ} \gu{કરે} \gu{છે}?} — \emph{śũ te tenī patnīne prem kare che?} — \textquotedblleft{}Does he love his wife?\textquotedblright{}
\end{quote}

\begin{cousinweb}
\textbf{\gu{પ્રેમ} \gu{કરવો}} (\emph{prem karvo}) is built from \emph{prem}, Sanskrit \emph{prēmán}, \textquotedblleft{}love\textquotedblright{}, taken over as a tatsama (CDIAL 8999; EtymDB), and the verb \emph{karvo}, \textquotedblleft{}to do\textquotedblright{}.
\end{cousinweb}

\begin{grammarlens}[title={What you've built}]
One more word for everyday life at A2.
\end{grammarlens}

\subsection*{Your turn}
\begin{itemize}
\raggedright
  \item \emph{Say it:} \emph{prem karvo}
  \item \emph{Say it:} \emph{prem karvo}, once more
  \item \emph{Say it:} \emph{śũ te tenī patnīne prem kare che?} — Does he love his wife?
  \item \emph{Recall:} say the Gujarati for to recognize, then the Gujarati for to embrace, then say \emph{prem karvo} again
\end{itemize}

\begin{tcolorbox}[breakable,colback=teal!4,colframe=teal!35!black,title={Before you move on}]
What does \gu{પ્રેમ} \gu{કરવો} mean? (\textquotedblleft{}To love\textquotedblright{}.) Say it once more.
\end{tcolorbox}
\section[\texorpdfstring{\gu{પરણવું} (paraṇvũ)}{પરણવું (paraṇvũ)}]{\gu{પરણવું} (paraṇvũ) — to marry, to get married}

\label{lesson:GU-C265-paranvu}



\begin{quote}
Before the new one: say the Gujarati for to embrace, then the Gujarati for to love.
\end{quote}

\subsection*{You'll want to know: \gu{પરણવું}}
\textbf{\gu{પરણવું}} — \emph{paraṇvũ} — \textquotedblleft{}to marry, to get married\textquotedblright{}.

\begin{quote}
\textbf{\gu{શું} \gu{તમારો} \gu{ભાઈ} \gu{આ} \gu{વરસે} \gu{પરણે} \gu{છે}?} — \emph{śũ tamāro bhāī ā varse paraṇe che?} — \textquotedblleft{}Is your brother getting married this year?\textquotedblright{}
\end{quote}

\begin{cousinweb}
\textbf{\gu{પરણવું}} (\emph{paraṇvũ}) is a tadbhava. It comes from Sanskrit \emph{páriṇayati}, \textquotedblleft{}leads round\textquotedblright{} the fire, and so \textquotedblleft{}marries\textquotedblright{} (CDIAL 7819).
\end{cousinweb}

\begin{grammarlens}[title={What you've built}]
One more word for everyday life at A2.
\end{grammarlens}

\subsection*{Your turn}
\begin{itemize}
\raggedright
  \item \emph{Say it:} \emph{paraṇvũ}
  \item \emph{Say it:} \emph{paraṇvũ}, once more
  \item \emph{Say it:} \emph{śũ tamāro bhāī ā varse paraṇe che?} — Is your brother getting married this year?
  \item \emph{Recall:} say the Gujarati for to embrace, then the Gujarati for to love, then say \emph{paraṇvũ} again
\end{itemize}

\begin{tcolorbox}[breakable,colback=teal!4,colframe=teal!35!black,title={Before you move on}]
What does \gu{પરણવું} mean? (\textquotedblleft{}To marry, to get married\textquotedblright{}.) Say it once more.
\end{tcolorbox}
\section[\texorpdfstring{\gu{સગાઈ} (sagāī)}{સગાઈ (sagāī)}]{\gu{સગાઈ} (sagāī) — engagement, betrothal}

\label{lesson:GU-C265-sagai}



\begin{quote}
Before the new one: say the Gujarati for to love, then the Gujarati for to marry.
\end{quote}

\subsection*{You'll want to know: \gu{સગાઈ}}
\textbf{\gu{સગાઈ}} — \emph{sagāī} — \textquotedblleft{}engagement, betrothal\textquotedblright{}.

\begin{quote}
\textbf{\gu{શું} \gu{તેની} \gu{સગાઈ} \gu{થઈ}?} — \emph{śũ tenī sagāī thaī?} — \textquotedblleft{}Did she get engaged?\textquotedblright{}
\end{quote}

\begin{cousinweb}
The origin of \textbf{\gu{સગાઈ}} (\emph{sagāī}) is not settled. It looks built on \emph{sagũ}, \textquotedblleft{}related\textquotedblright{}, with \emph{-āī}, but CDIAL has no entry for either word.
\end{cousinweb}

\begin{grammarlens}[title={What you've built}]
One more word for everyday life at A2.
\end{grammarlens}

\subsection*{Your turn}
\begin{itemize}
\raggedright
  \item \emph{Say it:} \emph{sagāī}
  \item \emph{Say it:} \emph{sagāī}, once more
  \item \emph{Say it:} \emph{śũ tenī sagāī thaī?} — Did she get engaged?
  \item \emph{Recall:} say the Gujarati for to love, then the Gujarati for to marry, then say \emph{sagāī} again
\end{itemize}

\begin{tcolorbox}[breakable,colback=teal!4,colframe=teal!35!black,title={Before you move on}]
What does \gu{સગાઈ} mean? (\textquotedblleft{}Engagement, betrothal\textquotedblright{}.) Say it once more.
\end{tcolorbox}
\section[\texorpdfstring{\gu{વહુ} (vahu)}{વહુ (vahu)}]{\gu{વહુ} (vahu) — daughter-in-law, bride}

\label{lesson:GU-C265-vahu}



\begin{quote}
Before the new one: say the Gujarati for to marry, then the Gujarati for engagement.
\end{quote}

\subsection*{You'll want to know: \gu{વહુ}}
\textbf{\gu{વહુ}} — \emph{vahu} — \textquotedblleft{}daughter-in-law, bride\textquotedblright{}.

\begin{quote}
\textbf{\gu{નવી} \gu{વહુ} \gu{ઘરમાં} \gu{નથી}.} — \emph{navī vahu gharmā̃ nathī} — \textquotedblleft{}The new bride is not in the house.\textquotedblright{}
\end{quote}

\begin{cousinweb}
\textbf{\gu{વહુ}} (\emph{vahu}) is a tadbhava. It comes from Sanskrit \emph{vadhū́}, \textquotedblleft{}bride, young wife\textquotedblright{}; CDIAL glosses the Gujarati word \textquotedblleft{}son's wife\textquotedblright{} (CDIAL 11250).
\end{cousinweb}

\begin{grammarlens}[title={What you've built}]
One more word for everyday life at A2.
\end{grammarlens}

\subsection*{Your turn}
\begin{itemize}
\raggedright
  \item \emph{Say it:} \emph{vahu}
  \item \emph{Say it:} \emph{vahu}, once more
  \item \emph{Say it:} \emph{navī vahu gharmā̃ nathī} — The new bride is not in the house.
  \item \emph{Recall:} say the Gujarati for to marry, then the Gujarati for engagement, then say \emph{vahu} again
\end{itemize}

\begin{tcolorbox}[breakable,colback=teal!4,colframe=teal!35!black,title={Before you move on}]
What does \gu{વહુ} mean? (\textquotedblleft{}Daughter-in-law, bride\textquotedblright{}.) Say it once more.
\end{tcolorbox}
\section[\texorpdfstring{\gu{વિશ્વાસ} (viśvās)}{વિશ્વાસ (viśvās)}]{\gu{વિશ્વાસ} (viśvās) — trust, faith}

\label{lesson:GU-C265-vishvas}



\begin{quote}
Before the new one: say the Gujarati for engagement, then the Gujarati for daughter-in-law.
\end{quote}

\subsection*{You'll want to know: \gu{વિશ્વાસ}}
\textbf{\gu{વિશ્વાસ}} — \emph{viśvās} — \textquotedblleft{}trust, faith\textquotedblright{}.

\begin{quote}
\textbf{\gu{શું} \gu{તમને} \gu{મારામાં} \gu{વિશ્વાસ} \gu{નથી}?} — \emph{śũ tamne mārāmā̃ viśvās nathī?} — \textquotedblleft{}Don't you trust me?\textquotedblright{}
\end{quote}

A masculine noun.

\begin{cousinweb}
\textbf{\gu{વિશ્વાસ}} (\emph{viśvās}) is Sanskrit \emph{viśvāsa}, \textquotedblleft{}trust, confidence\textquotedblright{}, taken over whole as a tatsama; CDIAL also lists an inherited Old Gujarati \emph{vīsāsa}, \textquotedblleft{}trust\textquotedblright{} (CDIAL 11966).
\end{cousinweb}

\begin{grammarlens}[title={What you've built}]
One more word for everyday life at A2.
\end{grammarlens}

\subsection*{Your turn}
\begin{itemize}
\raggedright
  \item \emph{Say it:} \emph{viśvās}
  \item \emph{Say it:} \emph{viśvās}, once more
  \item \emph{Say it:} \emph{śũ tamne mārāmā̃ viśvās nathī?} — Don't you trust me?
  \item \emph{Recall:} say the Gujarati for engagement, then the Gujarati for daughter-in-law, then say \emph{viśvās} again
\end{itemize}

\begin{tcolorbox}[breakable,colback=teal!4,colframe=teal!35!black,title={Before you move on}]
What does \gu{વિશ્વાસ} mean? (\textquotedblleft{}Trust, faith\textquotedblright{}.) Say it once more.
\end{tcolorbox}
